% Rubric: Discussion and Conclusions.
% Interpret findings in relation to the research question, draw
% evidence-supported conclusions, and identify limitations, uncertainty and
% potential improvements or applications without overstating the model scope.

\section{Discussion and conclusions}
\label{sec:discussion}

The tree topology explains RQ1's vulnerability: removing a station can cut the
only rail path between corridors, and extra capacity cannot restore a deleted
edge. RQ2 concerns whether surviving facilities acquire further overloads.
Dynamic containment reflects disappearing load, not preserved service; the
static all-station cascade is conditional on the forwarding rule and is not
evidence that a single closure would destroy Perth's railway.

RQ3 shows why recovery needs both service and facility metrics. Backup paths
can reconnect terminals while increasing load at particular rail facilities;
the Fremantle case traces that overload. Enough substitute connectivity can
maintain travel despite the extra closure, while a limited selection can leave
the network worse after secondary failures. Policies should
therefore be compared after the cascade stabilises, with a common intact OD
denominator and the reachable population reported beside travel times. Our
two-link budget counts selected connections, not fleet size.

Uncertainty remains broader than the Monte Carlo intervals in RQ4. We study
one service date and window, mapped undirected links, approximate rail speeds,
proxy capacities, symmetric effective bus times and immediate unlimited bus
standby. The model omits measured boarding/alighting, vehicle occupancy,
directional timetables, congestion, dispatch delay and finite substitute
capacity. Terminal-pair reachability is not proof of a feasible journey, and
the finite 86-station support limits tail inference: rejection cannot establish
the absence of scaling elsewhere.

Next tests would add observed OD demand and vehicle capacity, directional
time-dependent routing, delayed finite bus deployment, and several service-day
snapshots; enumerating all 86 triggers would remove trigger-sampling error. The
conclusions are narrower: targeted corridor removal is structurally damaging;
overload containment depends on the update rule and tolerance grid; and
substitute connectivity can improve service while aggravating rail overload.
These mechanisms are testable against operational evidence.
